\section{Preprocessing}
\label{sec:preprocessing}

Both datasets require transformation from raw event logs into the
$(n, d, A_{d,n})$ format consumed by our models, where $n$ indexes
users, $d$ indexes days within an experiment window, and $A_{d,n}$
is the trigger count for user~$n$ on day~$d$.

\subsection{REES46 Preprocessing}
\label{subsec:rees46_preprocessing}

\paragraph{Step 1: Chunked loading.}
The raw CSV files total approximately 47\,GB uncompressed.
To avoid memory exhaustion, we process each file in chunks of
2~million rows.
For each chunk, we extract the three relevant columns
(\texttt{event\_time}, \texttt{event\_type}, \texttt{user\_id}),
parse the timestamp, compute the calendar date, and aggregate
to user--day trigger counts immediately.
The per-chunk aggregations are then combined and re-aggregated
to handle users whose activity spans chunk boundaries.

This chunked approach reduces peak memory usage from
$\sim$100\,GB (loading all files at once) to $\sim$2--3\,GB,
making the pipeline feasible on standard hardware.

\paragraph{Step 2: User--day aggregation.}
After chunked loading, we obtain a table of
$(\texttt{user\_id}, \texttt{date}, \texttt{trigger\_count})$
triples.
We assign each date a zero-indexed \texttt{day\_index}
relative to the earliest date in the dataset and compute
each user's first trigger day.

\paragraph{Step 3: Summary statistics.}
The aggregated dataset exhibits the heavy-tailed engagement
pattern that motivates our methodology:
\begin{itemize}[leftmargin=*]
  \item The majority of users are active on very few days
        (median days active is low), while a small fraction
        exhibit sustained engagement over many days.
  \item Per-user total trigger counts follow a heavy-tailed
        distribution, consistent with the power-law behavior
        discussed in the main paper and documented broadly
        in online user activity data~\citep{clauset2009power}.
  \item Single-day users constitute a large fraction of the
        user base, underscoring the importance of modeling
        the ``long tail'' of low-engagement users.
\end{itemize}

\subsection{UCI Online Retail Preprocessing}
\label{subsec:uci_preprocessing}

\paragraph{Cleaning.}
Of the 541,909 raw rows, 135,080 (24.9\%) lack a customer
identifier (guest transactions) and are removed.
Cancellation invoices (prefixed with ``C'') and rows with
non-positive quantities are also removed, yielding 397,924 rows.

\paragraph{Aggregation.}
We define the daily trigger count $A_{d,n}$ for customer~$n$
on calendar day~$d$ as the number of distinct transaction line
items recorded for that customer on that day.
This produces 16,766 customer--day pairs across 4,339 unique
customers and 305 distinct calendar days.

\paragraph{User engagement characteristics.}
The UCI dataset exhibits qualitatively similar heavy-tailed
behavior:
35.7\% of customers are active on exactly one day,
67.4\% on three or fewer days,
while 0.6\% are active on more than 30~days.
The median total trigger count per customer is~41,
but the distribution has a long right tail
(maximum 7,847 triggers; standard deviation 228.8).

\subsection{Reproducibility}

All preprocessing scripts are included in the supplementary
code repository:
\begin{itemize}[leftmargin=*]
  \item \texttt{fitting/preprocess\_rees46.py} --- REES46 pipeline
        (chunked loading, aggregation, experiment construction);
  \item \texttt{lom\_revision/scripts/01\_preprocess\_uci.py} ---
        UCI Online Retail pipeline;
  \item \texttt{fitting/run\_full\_rees46\_pipeline.sh} ---
        end-to-end shell script (download, preprocess, fit, clean up).
\end{itemize}
The raw REES46 data is publicly available from
\url{https://data.rees46.com}; the UCI data from the
UCI Machine Learning Repository (ID~352).
Neither requires access credentials.
